\chapter{Introduction}

%\noindent\textbf{CHAPTER GUIDE:} Purpose: set context, define the problem, and state objectives + scope. Must include: Background, Problem statement, Objectives, Scope and limitations. How to write: mention domain + current approaches; quantify gaps with metrics. LaTeX: Cite with cite key, reference figures and tables. General rules: formal/technical style; cite non-trivial claims; reference figures/tables in text.

\section{Background}

% Test abbreviations and symbols
%This project uses \gls{ga} and \gls{ml} techniques \cite{deb2002nsga}. We also implement \gls{rl} with \gls{api} integration. The \gls{cpu} and \gls{gpu} work together for processing.

%Research shows optimization methods are crucial in modern computing \cite{smith2020}. Following IEEE guidelines \cite{ieee2022author, ieee2021editorial}, we implement state-of-the-art algorithms.

%The learning rate \gls{symb:alpha} is crucial, and we use discount factor \gls{symb:gamma}. Standard deviation \gls{symb:sigma} helps measure variance with mean \gls{symb:mu}.

Brain-Computer Interface (BCI) is a technology that allows a human to directly control a computer or other external device through activity in the brain. The idea of BCI has gained interest, not only because of its novelty, but also because of its practical significance to individuals who may be unable to interact with devices such as keyboards, mice, or touchscreens. So BCI has gained significance as a research area for its application in assistive communication, rehabilitation, and hands free computer interaction.

Electroencephalography (EEG) is commonly used for non invasive BCI systems as it records the electrical activity of the brain using electrodes placed on the scalp. It is the go to method because it is safe, low cost and portable but EEG signals are very small and easily affected by noise, so careful signal processing is necessary. 

Different BCI methods focus on certain EEG features, such as event related signals, visual responses or natural brain rhythms, to interpret the user’s intent such as alpha brain waves, which appear in the 8-12 Hz range and are usually linked to relaxed or unfocused states. As the person becomes attentive, there will be an increase in the beta wave, thus making it possible to use it as a control signal to determine the level of concentration. Eye blinks create signals that are very clear in the EEG recording. They can be used as control signals. Through the combination of the beta wave and the eye blinks, we will be able to develop a reliable BCI control mechanism without having to rely on classification methods. This project is an extension of that concept.
This system applies simple signal processing techniques like filtering and frequency analysis while avoiding or using few expensive equipment and advanced complex methods.

\section{Motivation}
The motivation for this project comes from the need for affordable and easy to use assistive technologies that support hands free communication for people in need. Our knowledge on signal filtering, processing and many other important elements can be implemented which enhances our capability. Beyond technology and skill it also serves in building and advancing our society by empowering the individuals with tool to communicate, interact and participate normally in regular life.

Another motivation is to explore whether BCI control can be achieved using basic EEG signals and low cost hardware. When eye blinks big potential spike is formed making the easiest EEG features to detect and understand compared to other brain signals. Using these signals along with simple processing techniques like filtering and Fast Fourier Transform (FFT), it is possible to build a BCI system that is easy to set up and use.

In addition, the motive behind the development of this project was to gain practical experience in EEG signal acquisition and processing. Development of a prototype of the BCI typing interface provided us with an understanding of the practical problems faced such as noise, constraints in hardware, economics, etc. As a result, this project has produced a prototype of the BCI system using very basic techniques in machine learning and hardware.

\section{Problem Statement}
Typing through regular computing relies on the use of physical components such as keyboards and mouses; however, these physical components cannot be utilized by individuals suffering from motor disabilities and situations that require more physical movement than what the individual can manage. Although brain-computer interface (BCI) solutions may provide a viable solution to this problem, most EEG-based brain-computer interface typing systems require expensive components and sophisticated machine learning algorithms that rely on huge amounts of data, thus making it impossible to use in low-budget settings.

The main problem that our project aims to address is the absence of a simple and affordable BCI typing system that works with simple EEG hardware and fundamental signal processing techniques. On the other hand, the important element to explore is whether eye blink detection and beta wave analysis can reliably control a basic typing or text selection interface without advanced classification algorithms and the difficult challenge is to extract usable control signals from noisy EEG data with few use of electrode while maintaining the system's reliability and usability. This project aims to solve this problem by designing and testing a low complexity EEG based BCI typing system suitable for educational and assistive purposes.


\section{Objectives}
The main objectives of this project are as follows:
\begin{itemize}
	\item To design a simple and low cost EEG based Brain-Computer Interface system
	\item To analyze EEG frequency bands using basic filtering and FFT
	\item To develop a basic text entry using detected EEG signals.
\end{itemize}



\section{Scope and Limitations}
This project focuses on designing and building a simple EEG based BCI typing or text selection system using eye blink detection and beta wave analysis via non invasive EEG signal recording, digital filtering, FFT based feature extraction and basic stimuli based decision making. Only a one EEG channel is used to make the system simple and cheap.

The project is not meant to attain fast typing speeds or be accurate on the same standards used commercially. There will be no incorporation of advanced machine learning and extensive testing for clinical validation. The project will serve the main purpose of educational and experimental demonstration of BCI principles and concepts.

The applications of this BCI System are:\begin{itemize}
	\item Assistive communication for people with limited motor abilities
	\item Hands free computer control where keyboards or mice are not practical
	\item Educational demonstrations for learning EEG signal processing and BCI fundamentals
	\item Student and research projects focusing on low cost neuroscience or biomedical engineering experiments
	\item Experimental platforms for further development of more advanced BCI.
\end{itemize}
